\documentclass[10pt]{article}
\usepackage[margin=0.9in]{geometry}
\usepackage{amsmath,amssymb,amsthm,microtype}
\usepackage[hidelinks]{hyperref}
\setlength{\emergencystretch}{3em}
\newtheorem{theorem}{Theorem}
\title{First Coefficient Contraction:\\A Sharper Seed Bound and a Lean Descent Certificate}
\author{Collatz proof project}
\date{September 5, 2026}
\begin{document}
\maketitle
\begin{abstract}
Before the multiplicative coefficient of a shortcut Collatz orbit first
contracts, its scaled correction is bounded linearly by its odd-step count.
This yields an explicit seed bound for a non-descending first contraction.
We combine that bound with two kernel-checked finite tables to prove that,
for every natural seed above one, any coefficient contraction within
65 steps has a prior or simultaneous descent. Every positive return of
length at most 65 consequently reaches one. No claim of mathematical
priority, global contraction, or a proof of Collatz is made.
\end{abstract}

\section{Exact correction and the first contraction}
Let $T(n)=n/2$ for even $n$ and $(3n+1)/2$ for odd $n$. Write
$x_t=T^t(N)$, let $j_t$ count odd values before time $t$, and define $d_t$ by
\[
2^t x_t=3^{j_t}N+d_t.
\]
The ideal correction is $c_t=d_t/3^{j_t}$. The coefficient contracts at
time $t$ when $3^{j_t}<2^t$. This is distinct from actual descent $x_t<N$:
the correction may offset a contracting coefficient.
\begin{theorem}\label{thm:correction}
If $2^s\le3^{j_s}$ for every $s<t$, then
\[
3d_t\le j_t3^{j_t},\qquad\text{equivalently}\qquad 3c_t\le j_t.
\]
\end{theorem}
\begin{proof}
Use induction on $t$. At an even state, both $j$ and $d$ remain unchanged.
At an odd state, $j_{t+1}=j_t+1$ and $d_{t+1}=3d_t+2^t$. The induction
hypothesis and the noncontraction condition at time $t$ give
\[
3d_{t+1}=9d_t+3\cdot2^t
 \le3j_t3^{j_t}+3\cdot3^{j_t}
 =(j_t+1)3^{j_t+1}.
\]
These forward recurrences follow from the existing correction cocycle.
\end{proof}
\begin{theorem}\label{thm:seed}
Under the same prefix hypothesis, if $3^{j_t}<2^t$ and $N\le x_t$, then
\[
3(2^t-3^{j_t})N\le j_t3^{j_t}.
\]
In particular, a seed strictly above
$j_t3^{j_t}/[3(2^t-3^{j_t})]$ descends at that endpoint.
\end{theorem}
\begin{proof}
The exact identity and $N\le x_t$ imply
$(2^t-3^{j_t})N\le d_t$. Apply Theorem~\ref{thm:correction}.
\end{proof}
The first-contraction hypothesis is essential to this proof. An arbitrary
paradoxical segment may have earlier coefficient contractions and is not
covered by this sharper correction estimate.

\section{A universal certificate through time 65}
\begin{theorem}\label{thm:finite}
For every natural $N>1$ and every $t\le65$,
\[
3^{j_t}<2^t\quad\Longrightarrow\quad
\exists s\le t\;\bigl(T^s(N)<N\bigr).
\]
\end{theorem}
First take the earliest coefficient contraction $r\le t$. The prefix
hypothesis of Theorem~\ref{thm:seed} then holds. A kernel-checked arithmetic
table proves, for all $r,j\in\{0,\ldots,65\}$ with $3^j<2^r$,
\[
j3^j<3(2^r-3^j)\cdot1186.
\]
Thus every seed $N\ge1186$ descends at its first contraction. The largest
integer seed bound in this table is 1185, at $(r,j)=(65,41)$.

For the remaining seeds $2\le N<1186$, a second kernel-checked table uses
a recursive Boolean checker. Its state is $(N,x,A,B,f)$, where $N$ is the
original seed, $x$ the current value, $A/B$ the current coefficient, and
$f$ the remaining horizon. It returns true immediately when $x<N$.
Otherwise it requires $B\le A$; if time remains, it advances $x$ once,
multiplies $A$ by three at an odd step, doubles $B$, and continues.
At zero remaining time it accepts exactly when $x<N$ or $B\le A$.

The soundness theorem states that acceptance implies, for every $u\le f$,
\[
A3^{j_u(x)}<B2^u\quad\Longrightarrow\quad
\exists v\le u\;\bigl(T^v(x)<N\bigr).
\]
Its proof is induction on the remaining horizon, with arbitrary $x,A,B$.
The finite table verifies acceptance from $(N,N,1,1,65)$ for every
$2\le N<1186$. It does not assume that every seed contracts within 65
steps: a non-descending orbit whose coefficient stays at least one through
the horizon is accepted too. This completes Theorem~\ref{thm:finite}.

\section{Returns and the remaining gap}
A positive return $T^q(N)=N$ with $q>0$ has $3^{j_q}<2^q$ by the existing
cycle inequality. If $q\le65$ and $N>1$, Theorem~\ref{thm:finite} supplies
a smaller positive state on that return orbit. It has the same $q$-step
return. Strong induction on the seed therefore proves
\[
0<N,\quad 0<q\le65,\quad T^q(N)=N
\quad\Longrightarrow\quad\exists s\;T^s(N)=1.
\]
This finite-horizon result is separate from the Python census of arbitrary
paradoxical segments at length 65. The latter allows earlier descent and
found 244 segments. The new theorem verifies that such descent must occur
for every seed above one whenever the length-65 coefficient contracts;
it does not kernel-certify the census's exact list.

The general seed bound holds at all first contraction times, but the two
finite tables cover only 65 steps. Neither existence of a coefficient
contraction for every seed nor the required descent implication at all
later times is proved. Unbounded orbits and longer nontrivial cycles remain
unresolved.

\section{Formal verification}
{\raggedright
Source: \texttt{lean/Collatz/FirstContraction.lean}. Principal declarations
are \texttt{first\_contraction\_correction\_bound},
\texttt{first\_contraction\_seed\_bound},
\texttt{contractionCheck\_sound},
\texttt{descent\_of\_contraction\_le65}, and
\texttt{return\_le65\_reaches\_one}. Both finite tables use ordinary
\texttt{decide}, so their proof terms are checked by kernel reduction;
no \texttt{native\_decide}, custom axiom, or \texttt{sorry} is used.
The module build and selected 131-declaration axiom audit pass using only
standard foundations. The audit inspects dependencies; it is not a second
independent kernel implementation. Build with
\texttt{lake build Collatz.FirstContraction} from \texttt{lean/}.\par}
\end{document}
